% Auto-generated --- do not edit by hand
\documentclass[a4paper,10pt]{article}

\usepackage[T1]{fontenc}
\usepackage[utf8]{inputenc}
\usepackage{tgheros}
\renewcommand{\familydefault}{\sfdefault}
\usepackage[top=36pt,bottom=34pt,left=50pt,right=50pt]{geometry}
\usepackage[dvipsnames]{xcolor}
\usepackage{graphicx}
\usepackage{tikz}
\usepackage[hidelinks]{hyperref}
\usepackage{fontawesome5}
\usepackage{parskip}
\usepackage{tabularx}
\usepackage{calc}
\usepackage{microtype}

% ── Colors ──
\definecolor{cvblack}{RGB}{35,35,40}
\definecolor{cvdark}{RGB}{50,50,58}
\definecolor{cvgray}{RGB}{110,110,118}
\definecolor{cvlight}{RGB}{165,165,172}
\definecolor{accent}{RGB}{160,125,80}
\definecolor{ruledark}{RGB}{60,60,68}
\definecolor{rulelight}{RGB}{200,200,206}

\color{cvblack}

% ── Layout ──
\pagestyle{empty}
\setlength{\parindent}{0pt}
\setlength{\parskip}{0pt}

\newlength{\datecol}
\setlength{\datecol}{88pt}
\newlength{\colgap}
\setlength{\colgap}{10pt}

% ── Icons ──
\newcommand{\iconlabel}[2]{%
  {\fontsize{7.5}{9}\selectfont\color{cvgray}{\fontsize{8}{9}\selectfont#1}\enspace#2}%
}

% ── Section heading ──
\newcommand{\cvsection}[1]{%
  \vspace{6pt}%
  {\fontsize{9}{11}\selectfont\bfseries\color{cvdark}\MakeUppercase{#1}}%
  \\[-3pt]%
  {\color{ruledark}\rule{\textwidth}{0.6pt}}%
  \vspace{4pt}%
}

% ── Experience entry ──
\newcommand{\cventry}[5]{%
  \noindent
  \begin{tabularx}{\textwidth}{@{}p{\datecol}@{\hspace{\colgap}}X@{}}
    {\fontsize{8}{10}\selectfont\color{cvgray}#1} &
    {\fontsize{9.5}{12}\selectfont\bfseries\color{cvblack}#2}\newline
    {\fontsize{8.5}{11}\selectfont\color{cvgray}#3}\newline
    {\fontsize{8.5}{11}\selectfont\color{cvblack}#4}\newline
    {\fontsize{8}{10}\selectfont\color{accent}#5}
  \end{tabularx}\vspace{3pt}%
}

% ── Education entry ──
\newcommand{\cvedu}[4]{%
  \noindent
  \begin{tabularx}{\textwidth}{@{}p{\datecol}@{\hspace{\colgap}}X@{}}
    {\fontsize{8}{10}\selectfont\color{cvgray}#1} &
    {\fontsize{9.5}{12}\selectfont\bfseries\color{cvblack}#2}\newline
    {\fontsize{8.5}{11}\selectfont\color{cvgray}#3}\newline
    {\fontsize{8.5}{11}\selectfont\color{cvblack}#4}
  \end{tabularx}\vspace{3pt}%
}

% ── Skill entry ──
\newcommand{\cvskill}[2]{%
  {\fontsize{9}{12}\selectfont{\bfseries\color{cvblack}#1:\enspace}{\color{cvgray}#2}}\\[3pt]%
}

% ── Talk entry ──
\newcommand{\cvtalk}[2]{%
  {\fontsize{9}{12}\selectfont{\bfseries\color{cvblack}#1\enspace}{\color{cvgray}#2}}\\[3pt]%
}

\begin{document}

% ── Header ──
\noindent
\begin{minipage}[t]{\textwidth-82pt}
  {\fontsize{22}{26}\selectfont\bfseries\color{cvblack}Thomas Rolf Mannhart}\\[3pt]
  {\fontsize{11}{14}\selectfont\color{accent}Professional AI Engineer}\\[5pt]
  \begin{tabular}{@{}l@{\hspace{18pt}}l@{}}
    \iconlabel{\faIcon{map-marker-alt}}{Zürich, Switzerland} &
    \iconlabel{\faEnvelope}{\href{mailto:thomas@mannhart.ai}{thomas@mannhart.ai}} \\[3pt]
    \iconlabel{\faGithub}{\href{https://github.com/thommann}{github.com/thommann}} &
    \iconlabel{\faGlobe}{\href{https://t.mannhart.ai}{t.mannhart.ai}}
  \end{tabular}
\end{minipage}%
\hfill
\begin{minipage}[t]{72pt}
  \vspace{0pt}
  \begin{tikzpicture}
    \clip (0,0) circle (36pt);
    \node at (0,0) {\includegraphics[width=72pt]{/home/thomas/Projects/mannhart.ai/src/assets/img/photo-cv.jpg}};
  \end{tikzpicture}
\end{minipage}

\vspace{4pt}
{\color{ruledark}\rule{\textwidth}{0.8pt}}
\vspace{2pt}

% ── Profile ──
\cvsection{Profile}
{\fontsize{9}{12}\selectfont\color{cvblack}I'm a Professional AI Engineer at bbv Software Services in Zürich. Together with my team, I'm building the Swiss AI Hub --- a Swiss-made, model-agnostic AI platform for organizations of any size. Beyond the platform, I design and implement custom AI solutions for our customers, with a focus on RAG systems and agentic workflows. I also serve as technical lead on customer AI projects and consult organizations on their AI strategy.

\vspace{4pt}
I studied Informatics at the University of Zürich. I did my Bachelor's in Software Systems and my Master's in Artificial Intelligence. Already during my Bachelor's, I started working as a software engineer. My first role was at swissbiomechanics, an ETH spin-off, where I independently built a Java application for tracking biomedical analyses. At PolygonSoftware, a startup founded by UZH graduates, I led full-stack development and worked on computer vision and machine learning projects. After finishing my Master's, I joined Ergon Informatik, one of Zürich's most respected software houses, where I worked on enterprise software for the retail sector. From there, I moved to bbv, which offered me the position of AI Engineer.}

% ── Experience ──
\cvsection{Experience}
\cventry{2025--now}{Professional AI Engineer}{bbv Software Services AG, Zürich}{Development of a comprehensive enterprise AI platform, including architecture design in collaboration with the software architect. Design and implementation of customized AI solutions --- especially RAG systems --- for customer projects in industry and market research. Technical leadership (Dev Lead) of customer AI projects. Operation and maintenance of the AI platform and customer-specific solutions. Consulting customers on their IT and AI strategy.}{Python, LLMs / RAG, Agentic AI, Dev Lead, AI Strategy, Azure, TypeScript, Platform Engineering}
\cventry{2023--2024}{Professional Software Engineer}{Ergon Informatik AG, Zürich}{Developed a time-tracking and workforce planning system for the retail sector. Handled end-to-end software delivery from requirements engineering and prototyping through to second/third-level support. Mentored new team members and contributed to youth outreach by organizing IT workshops for students.}{Java, Kotlin, Angular, TypeScript, SQL, Selenium, Jenkins, Git}
\cventry{2020--2023}{Senior Software Developer}{PolygonSoftware, Opfikon}{Led full-stack development of web applications and computer vision / machine learning projects at a UZH-founded startup. Projects included an OCR pipeline for a logistics company and a data visualization application for vibration measurement devices. Designed software architectures, supervised dev teams, and interfaced directly with product owners and clients.}{Full Stack, Computer Vision, Machine Learning, Web Apps, DevOps}
\cventry{2019--2020}{Junior Software Developer}{swissbiomechanics ag, Zürich}{Led an independent software project for this ETH spin-off, building a Java application to track biomedical analyses and automatically generate clinical reports. Handled stakeholder communication, requirements analysis, and coordinated with other developers.}{Java, Report Generation, Biomedical}

\newpage

% ── Education ──
\cvsection{Education}
\cvedu{2020--2023}{MSc in Informatics (Artificial Intelligence)}{University of Zürich}{Specialization in Artificial Intelligence. Exchange semester at the National Taiwan University (NTU) in Taipei. Team project (with Johann Schwabe and Joel Barmettler): Airspace Auction Simulator. Thesis: KroneDB --- Compressing and Querying Time Series Data using the Kronecker Decomposition. Supervised by Johannes Marti and Dan Olteanu at the Data Systems and Theory group.}
\cvedu{2017--2020}{BSc in Informatics (Software Systems)}{University of Zürich}{Specialization in Software Systems. Thesis: A General-purpose Range Join Algorithm for PostgreSQL. Written at the Free University of Bozen-Bolzano under the guidance of Anton Dignös. Awarded the Semesterpreis HS20 by the Faculty of Business, Economics and Informatics --- one of four students recognized across the entire faculty that semester.\newline{\bfseries\color{accent}\href{https://www.ifi.uzh.ch/en/archive/news-archive/outstanding-scientific-work-thomas-mannhart.html}{Semesterpreis HS20}}}

% ── Skills ──
\cvsection{Skills}
\cvskill{Programming}{Python (expert), TypeScript / JS (advanced), Java / Kotlin (advanced), SQL (advanced)}
\cvskill{AI}{Agent Orchestration (expert), RAG (expert), MCP (advanced)}
\cvskill{Tools}{Claude Code (expert), Git / GitHub (expert), Docker (advanced)}
\cvskill{Languages}{German (native), English (fluent), French (basic)}

% ── Talks ──
\cvsection{Talks \& Projects}
\cvtalk{KI als Entwicklungspartner}{bbv webinar on practical methods, tools, and proven strategies for integrating AI into the software development lifecycle --- from requirements analysis through code development to testing.}
\cvtalk{AI-Augmented Software Engineering}{Talk at the FHNW Data Science \& Data Engineering Alumni Event on integrating AI into the software development lifecycle --- from enterprise AI strategy and the Swiss AI Hub, to working with coding agents and the "Development Funnel" for context-driven AI workflows.}
\cvtalk{About the Swiss AI Hub}{Swiss-made, model-agnostic enterprise AI platform. Listed on Siemens Xcelerator, certified Swiss Made Software.}

\end{document}
